\section{Completion from a separated boundary}
\label{sec:boundary-completion}

We now complete the pairing from the separation stop.  The main
constraint is to control individual resolvent entries as the unresolved
set shrinks.  We compare the resolvent with a smooth path sum fixed at the
stop and continued by exact matrix updates.  Separation makes its rows
nearly single coordinates.  An initial averaged energy gap then gives
small individual row energies, which in turn permit entry concentration
down to a small final matching.
The stop configuration and the completed matching will define one reference
throughout the primary parameter interval.  Keeping this same reference as
$h$ varies is essential for the local spectral estimates in the next section.

The inputs from the preceding sections are the following:
\begin{enumerate}
\item The $l_0$ unresolved stubs lie at distinct vertices.  No nonempty
actual nonbacktracking path of length at most $R_0$ joins two of these
vertices, including a return to the same vertex.
\item At every retained primary and both signs, the conclusions of
Proposition~\ref{prop:cleanup} and
Lemma~\ref{lemma:local-laws-implications} hold.  In particular, the
precision is positive; the trace, diagonal-square, entry and local-energy
bounds hold; and $\mathcal T$ has the stated lower bound.
\item The coefficient and actual-path multiplicity bounds through length
$L$ hold, including their versions with any fixed power of the path
length.
\end{enumerate}
Only the primary parameters are retained, with both signs.  A uniform
pair means a uniformly chosen ordered pair of distinct unresolved stubs;
reversing its order gives the same next state.  These finite averages are
used only to implement the deterministic account rule.  Vertex operators
act on $\mathbf1^\perp$ and are extended by zero when entries are taken.
Constants are uniform in $n$ and the retained points, but may depend on
the fixed ratio bounds for $L/R_0$.  Uniformity in $h$ always refers to the
convex hull of the retained primaries, contained in $(0,h_*]$.
The stop precision is positive throughout this interval for both signs:
the identity $-\partial_hM=\mathcal D+h^{-2}\alpha I\succ0$ makes
$M(h)\succeq M(h_*)\succ0$ when $h\le h_*$.

Separation makes every subsequent pair of distinct unresolved vertices an
allowable simple edge.  Indeed, if a path of length at most $R_0$ from one
unresolved vertex used a newly added edge, the segment preceding the first
new edge would be an actual stop-graph path to an old boundary vertex.
That segment has positive length unless it starts at the unresolved
vertex itself; the latter vertex is not incident to any previously
added edge.  Separation excludes both a short path to a different
unresolved vertex and a short return.  This argument applies at every
subsequent state and for every matching of the boundary.  Thus the
uniform pair averages below require no further graph-theoretic rejection.

\begin{samepage}
\begin{proposition}[Completion of the separated boundary]
\label{prop:boundary-completion}
Under these inputs, the account rule can deterministically complete the
pairing.  Every added edge is simple and preserves the girth exclusion.
At the end the precision is positive at every primary point for both
signs.  There is a reference matrix $\widehat P_{\rm fin}$, defined below
by exact numerical updates from a smooth coefficient truncation at the
stop, such that at every retained primary parameter and both signs,
\[
 \max_{i,j}|(W_{\rm fin}-\widehat P_{\rm fin})_{ij}|
 \le C B_B,
 \qquad B_B=\frac{n^{.40a}}{\sqrt{l_2}}.
\]
The construction retains the stop configuration and the matching added
after it.  These define the same family $\widehat P_{\rm fin}(h)$ throughout
the primary interval.
All the accounts, candidate values and transitions in this section can
be evaluated in a number of bit operations polynomial in $n$, with an
exponent depending only on the fixed parameters.
\end{proposition}
\end{samepage}

\subsection{The reference at the separation stop}
\label{subsec:boundary-reference}

The reference must satisfy two different estimates: most of each row on
the boundary must be small, and its compression to boundary contrasts
must be close to the identity in operator norm.  The first follows from
separation and path counting.  The second needs the smooth cutoff below.

Let $S_0$ be the unresolved set at the stop and let
$\mathcal H_0=\mathbf 1_{S_0}^{\perp}$.  Sending a stub coordinate to its
vertex coordinate gives an isometry
$V_0:\mathcal H_0\longrightarrow\mathbf 1^\perp$.  Thus
$\mathcal D=V_0V_0^T$.  We use matrices on $\mathcal H_0$ also as
zero-extended matrices in $S_0$ coordinates.  In particular its identity
has coordinate matrix $\Pi_{l_0}$.

Fix an even, compactly supported function $\varphi$ such that
$\varphi(x)=1$ for $|x|\le 1-\delta_L$ and $\varphi(x)=0$ for
$|x|\ge1$, where $0<\delta_L\le .03a$.  Its finite smoothness order
will be chosen in the proof below.  A piecewise polynomial with rational
coefficients, and with the specified derivatives vanishing at its
joins, suffices.  Define
\[
 \widehat P_0=\sum_{k\ge0}\varphi(k/L)(\sigma h)^k Q_k
 \quad\hbox{on }\mathbf 1^\perp.
\]
The sum is finite.  By the coefficient estimates,
\[
 \max_{i,j}|(\widehat P_0-P_{<L})_{ij}|
 \le \frac{n^{.32a}}{\sqrt{l_0}}.
\]
Indeed the altered actual terms have lengths at least
$(1-\delta_L)L$, and hence their absolute sum at fixed endpoints is
$O(Lq^{-(1-\delta_L)L/2})$; the virtual and centering contributions are
$O(n^\eta q^{L/2}/l_0)$.  Substitution of the prescribed value of $L$
gives the displayed bound.  The same argument applies to centered-stub
outputs and to a fixed number of $h$ derivatives, allowing an
arbitrarily small additional power loss.

\begin{lemma}[Sparse rows and a small compressed operator]
\label{lem:boundary-smooth-reference}
There are constants $r_P'>0$ and $C$ such that
\begin{equation}
 \widehat\Gamma_0:=V_0^T\widehat P_0V_0
   =I+\widehat D_0\quad\hbox{on }\mathcal H_0,
 \qquad \|\widehat D_0\|\le \frac{C}{R_0}.
 \label{eq:original-40}
\end{equation}
For every global vertex coordinate $i$, the vector
$p_{i,0}=V_0^T\widehat P_0e_i$ satisfies
\begin{equation}
 p_{i,0}=c_i\Pi_{l_0}e_{s(i)}+\tau_i,
 \qquad |c_i|\le1,\qquad \|\tau_i\|_2\le n^{-r_P'}.
 \label{eq:original-41}
\end{equation}
The first summand is omitted if there is no actual path from $i$ to a
boundary vertex of length less than $\lfloor R_0/4\rfloor$.  Otherwise
that vertex and that path are unique, and
$c_i=(\sigma h)^{\operatorname{dist}(i,s(i))}$.
For a boundary vertex, $s(i)$ is its own stub and $c_i=1$.
The estimate for $\tau_i$ holds also for $\partial_h\tau_i$, after
decreasing $r_P'$ if necessary.  All these statements hold uniformly
throughout the closed convex hull of the retained primary parameters.
\end{lemma}

\begin{proof}
\textit{Short paths and row tails.}
Put $t_0=\lfloor R_0/4\rfloor$.  Two boundary endpoints within distance
$t_0$ of a given vertex would give a boundary-to-boundary path of length
at most $2t_0$.  Two distinct nonbacktracking paths of length less than
$t_0$ from a boundary endpoint to a fixed vertex give, after cancellation
of a common terminal segment, a nonempty nonbacktracking boundary
return of length less than $2t_0$.  Both possibilities are excluded.
This proves uniqueness of the short summand in
\eqref{eq:original-41}.

For $t_0\le k<L$, fix a path prefix of length $k-t_0$ starting at $i$.
There is at most one boundary endpoint within distance $t_0$ of its
last vertex.  The bounded fixed-endpoint multiplicity estimate gives
at most a constant number of nonbacktracking suffixes of length $t_0$
to that endpoint.  Consequently there are at most $Cq^{k-t_0}$ actual
paths of length $k$ from $i$ to the boundary.  Applying the same
fixed-endpoint multiplicity estimate to the complete paths yields
\[
 \sum_{s\in S_0}
 \left|\sum_{t_0\le k<L}\varphi(k/L)(\sigma h)^k
           Q_{k,\mathrm{actual}}(i,s)\right|^2
 \le C L\sum_{t_0\le k<L}q^{-k}q^{k-t_0}
 \le n^\eta q^{-t_0}.
\]
The virtual and centering contribution to this row has norm at most
$C(1+L)^2q^{L/2}/\sqrt{l_0}=C(1+L)^2n^{-.29a}$.
Projection onto $\mathcal H_0$ is a contraction.  These bounds prove
\eqref{eq:original-41}; insertion of factors $k/h$ proves the derivative
statement.  The same argument for a boundary input, and symmetry,
gives
\[
 \max_{s\in S_0}\|\widehat D_0\Pi_{l_0}e_s\|_2
 +\max_{s\in S_0}\|\partial_h\widehat D_0\Pi_{l_0}e_s\|_2
 \le n^{-r_P'}
\]
with a possibly smaller exponent.  This completes the row estimates.

\par\smallskip\noindent
\textit{Moments on boundary contrasts.}
For the operator estimate, we use the stronger consequence of separation
$V_0^TQ_jV_0=0$ for $1\le j\le R_0$.
To verify this also for paths with virtual traversals, observe that
the first virtual traversal must be the first step: a positive actual
prefix would already join a boundary vertex to a missing departure.
After a virtual traversal the sole missing stub at its landing vertex
is excluded as the next departure.  Any later virtual traversal, or
an actual endpoint on the boundary, would therefore require a positive
actual boundary-to-boundary segment of prohibited length.  The only
remaining possibility is the single virtual traversal; its averaging
operator vanishes on $\mathcal H_0$.

To express this vanishing in terms of spectral moments, introduce, for a
formal variable $t$, the matrix series
\begin{align*}
 G_{\rm gen}(t)
  &=V_0^T(1-t^2)
    [I-tA_e+t^2(qI-V_0V_0^T)]^{-1}V_0,\\
 G_{\rm reg}(t)
  &=V_0^T(1-t^2)[I-tA_e+qt^2I]^{-1}V_0.
\end{align*}
The formal Woodbury identity and the preceding vanishing give
\[
 G_{\rm gen}(t)=I+O(t^{R_0+1}),\qquad
 G_{\rm reg}(t)
 =G_{\rm gen}(t)
    \left(I+\frac{t^2}{1-t^2}G_{\rm gen}(t)\right)^{-1}
 =(1-t^2)I+O(t^{R_0+1}).
\]
Let $\mathcal U_k$ denote the Chebyshev polynomial of the second kind,
normalized by
$\mathcal U_k(\cos\vartheta)
 =\sin((k+1)\vartheta)/\sin\vartheta$.
For every complex unit vector $f\in\mathcal H_0$, the spectral measure
$\mu_f$ of $A_e/(2\sqrt q)$ at $V_0f$ therefore satisfies
\[
 \int\mathcal U_0\,d\mu_f=1,
 \qquad \int\mathcal U_k\,d\mu_f=0
       \quad(1\le k\le R_0).
\]
This follows on dividing by $1-t^2$ and using
$(1-2xy+x^2)^{-1}=\sum_{k\ge0}\mathcal U_k(y)x^k$ with
$x=t\sqrt q$.  Hence $\mu_f$ agrees, through polynomial degree $R_0$,
with the semicircle measure of angular density
$(2/\pi)\sin^2\vartheta\,d\vartheta$.

Positivity at the two closest signed precisions implies
$\|A_e|_{\mathbf1^\perp}\|\le2\sqrt q+C\alpha_*^2$.
For example, multiplying the positive precision by $h$ and adding the
positive term $h^2V_0V_0^T$ gives
$I-\sigma hA_e+qh^2I\succ0$.
The claimed support bound follows from
$z(h)=2\sqrt q+O(\alpha_*^2)$.

\par\smallskip\noindent
\textit{Edge mass and the smooth cutoff.}
The moment identities give the following bound: for every
$u\ge C'/R_0$, the mass within angular distance $u$ of either edge,
including mass outside $[-1,1]$ at that edge, is at most $Cu^3$.
For its proof, choose $m=\lfloor c/u\rfloor$ with a sufficiently small
positive constant $c$, so that $2m\le R_0$, and put
$K_m(y)=\sum_{k=0}^m(k+1)\mathcal U_k(y)$.
Moment matching and orthonormality give
\[
 \int K_m(y)^2\,d\mu_f(y)
 =\sum_{k=0}^m(k+1)^2\le C(m+1)^3.
\]
For $0\le\vartheta\le c_0/(m+1)$ every summand satisfies
$\mathcal U_k(\cos\vartheta)\ge c_1(k+1)$; hence
$K_m\ge c_2(m+1)^3$ there.  For $y\ge1$ the same inequality follows
from $\mathcal U_k(\cosh v)=\sinh((k+1)v)/\sinh v\ge k+1$.
Markov's inequality gives mass $O((m+1)^{-3})=O(u^3)$.
Reflection proves the estimate at $-1$.  Values of $u$ bounded away
from zero are covered by increasing the constant.

Put $x=t\sqrt q$ and write
$H_b(x)=G_{\rm reg}(t)-(1-t^2)I=\sum B_kx^k$.
Its coefficients vanish for $k\le R_0$.
Choose a smooth function $\psi(k/L)$ supported in
$R_0/2<k<C L$ and equal to one for $R_0<k\le L$.
The ratio $L/R_0$ is bounded above and below by fixed positive
constants, so all derivatives of this cutoff, in the variable $k/L$,
have fixed bounds.  For $x_0=\sigma h\sqrt q$ and $\xi\in\mathbb R$,
we claim
\begin{equation}
 \left\|\sum_{k>R_0}B_kx_0^k\psi(k/L)e^{i\xi k/L}\right\|
 \le \frac{C(1+|\xi|)^{J_0}}{R_0}
 \label{eq:original-42}
\end{equation}
for a fixed integer $J_0$.

For $k>R_0$, $B_k$ is the compression to $\mathcal H_0$ of
$\mathcal U_k(A_e/(2\sqrt q))-q^{-1}\mathcal U_{k-2}(A_e/(2\sqrt q))$.
Thus each quadratic form in \eqref{eq:original-42} is the integral of
a tapered Chebyshev kernel against $\mu_f$.  The cubic edge-mass bound
will control the kernel's $O(L^2)$ edge peaks, up to a fixed power of
$1+|\xi|$; smoothness of the cutoff will give summable decay away
from the edges.  We now establish these kernel bounds.

Extend the coefficient
sequence by zero to all integers.  For
$b_k=\psi(k/L)|x_0|^ke^{i\xi k/L}$ and every fixed integer $J\ge0$,
the $J$th finite difference has $\ell^1$ norm at most
$C_JL^{1-J}(1+|\xi|)^J$.
This follows by the fundamental theorem of calculus applied $J$
times to the smooth interpolation of $b_k$; derivatives of
$|x_0|^k$ are harmless because, on the support $k\ge cL$,
$(L|\log|x_0||)^j e^{-cL|\log|x_0||}$ is bounded.
Summation by parts, or the trivial $\ell^1$ bound when the denominator
is small, yields
\[
 \left|\sum_kb_ke^{ik\omega}\right|
 \le C_J L(1+|\xi|)^J
   (1+L\operatorname{dist}(\omega,2\pi\mathbb Z))^{-J}.
\]
The sine formula for the two Chebyshev terms, the preceding
bound, and $|\mathcal U_k(\cos\vartheta)|\le k+1$ imply, after
increasing $J$ if necessary,
\[
 |\mathcal K(\cos\vartheta)|
 \le C_JL^2(1+|\xi|)^J
   (1+L\min(\vartheta,\pi-\vartheta))^{-J}.
\]
If $y$ lies just outside $[-1,1]$, its hyperbolic angle is
$O(\alpha_*)$; since $L\alpha_*=o(1)$, the corresponding bound is
$C_JL^2(1+|\xi|)^J$.

Integrate first over the edge bins of angular width $C'/R_0$.
Their mass is $O(R_0^{-3})$, so their contribution is
$O(L^2R_0^{-3})(1+|\xi|)^J=O(R_0^{-1})(1+|\xi|)^J$.
The subsequent bins, with angular distance between
$2^jC'/R_0$ and $2^{j+1}C'/R_0$, contribute at most
\[
 C L^2(1+|\xi|)^J
       (2^{j+1}/R_0)^3(1+c2^jL/R_0)^{-J}.
\]
For $J>3$ their sum is again $O(R_0^{-1})(1+|\xi|)^J$.
This holds for every complex unit $f$.  The numerical-radius
inequality $\|T\|\le2\sup_{\|f\|=1}|f^*Tf|$ proves
\eqref{eq:original-42}.  If $x_0<0$, replace $y$ by $-y$ in this
argument; the two edge bounds are symmetric.

\par\smallskip\noindent
\textit{Return to the configuration pencil.}
The preceding bound concerns the regular pencil.  To recover the desired
smooth truncation of the configuration pencil, eliminate in the opposite
direction:
\[
 G_{\rm gen}(t)
 =I+\sum_{j\ge1}\frac{t^{2(j-1)}}{(1-t^2)^{2j}}H_b(x)^j.
\]
Only $j\le L/(R_0+1)$ can contribute below total degree $L$.
Choose once and for all an integer $J_*$ larger than all these upper
bounds.  Insert $\psi(k/L)$ into each factor $H_b$: this does not
change any coefficient contributing after the final taper
$\varphi(k/L)$.  Fourier inversion expresses that taper as
\[
 \varphi((k_1+\cdots+k_j+m)/L)
 =\frac1{2\pi}\int_{\mathbb R}\widehat\varphi(\xi)
             e^{i\xi(k_1+\cdots+k_j+m)/L}\,d\xi.
\]
At $|t|\le1/\sqrt q$ the scalar series
$t^{2(j-1)}(1-t^2)^{-2j}$ has a bounded absolute sum for each
$j\le J_*$.  The $j$ matrix factors inside the integral are bounded
by \eqref{eq:original-42}.  Choose the finite smoothness of
$\varphi$ so that
$\int |\widehat\varphi(\xi)|(1+|\xi|)^{J_0J_*}\,d\xi<\infty$.
The tapered $j$th term then has norm $O(R_0^{-j})$.
Summing the finitely many terms proves
\eqref{eq:original-40}.
\end{proof}

\subsection{Continuing the reference through every matching}

We continue the reference by the same numerical update as the resolvent.
An aggregate formula for all pairs added since the stop will show that
its estimates persist for every possible matching.

Let $P_t$ be the stub operator on $S_0$ consisting of the resolved
swaps and uniform averaging on the currently unresolved stubs, and
restrict it to $\mathcal H_0$.  Let $Q_R$ be the orthogonal projection
onto the current unresolved contrasts in this fixed space.  Then
$P_t^2=I-Q_R$, $P_tQ_R=Q_RP_t=0$, and the precision increment since
the stop is
\[
 V_0\,[h(I-Q_R)-\sigma P_t]V_0^T.
\]
The Woodbury coefficient at ideal compressed Gram $I$ is
$H_{t0}=-\sigma hP_t$.  This identity is checked on the three
eigenspaces of $P_t$, with eigenvalues $0,1,-1$, using
$m=h/(1-h^2)$.  Define
\[
 H_t=H_{t0}(I+\widehat D_0H_{t0})^{-1},\qquad
 \widehat P_t=\widehat P_0
       -\widehat P_0V_0H_tV_0^T\widehat P_0.
\]
By \eqref{eq:original-40}, the inverse exists and $\|H_t\|\le C$.
The identity
$H_t=(I+H_{t0}\widehat D_0)^{-1}H_{t0}$ also shows that
$H_tQ_R=Q_RH_t=0$.  Thus these are precisely the iterated exact
numerical updates of the reference; no positivity of the reference
is asserted or needed.  The formula depends only on the stop data and
the current matching, so it defines the reference simultaneously for
all $h$ in the primary interval, irrespective of the order in which
the matching was exposed.

On current contrasts its compressed error is
\[
 \widehat D
 =Q_R(\widehat D_0-\widehat D_0H_t\widehat D_0)Q_R,
 \qquad \|\widehat D\|\le C/R_0.
\]
The row norms of this matrix are power-small, uniformly over every
matching of $S_0$.  To see this without accumulating a stepwise
error, note that $P_t\Pi_{l_0}e_s$ is a partner coordinate or an
average of coordinates, followed by the fixed centering.  Therefore
\eqref{eq:original-41} gives
$\|\widehat D_0H_{t0}\Pi_{l_0}e_s\|\le Cn^{-r_P'}$.
The identity
\[
 \widehat D_0H_t
 =(I+\widehat D_0H_{t0})^{-1}\widehat D_0H_{t0}
\]
preserves this bound.  Projection to current contrasts replaces a
coordinate by that coordinate minus an average, so the same estimate
holds for rows of $\widehat D$.

For every global $i$, its row on current unresolved stubs is
\begin{equation}
 p_i=c_i\Pi_le_{s(i)}+\tau_{i,t},\qquad
 \|\tau_{i,t}\|_2\le n^{-r_P},
 \label{eq:original-43}
\end{equation}
if $s(i)$ is still unresolved, and consists only of the stated tail
otherwise.  Indeed the row is
$Q_R[p_{i,0}-\widehat D_0H_tp_{i,0}]$, and the preceding bounds
apply separately to its short coordinate and its tail.
We may fix $r_P=.001a$, decreasing the intermediate exponent if
needed.  Initially the diagonal estimates from the path bounds also
give
\[
 |\operatorname{Tr}\widehat D_0|\le n^\eta q^{L/2},\qquad
 \sum_s(\widehat D_0)_{ss}^2\le n^\eta.
\]

\subsection{Pair updates and the initial energy gap}

We now turn to the actual resolvent.  Its compressed deviation from the
identity enters each pair update.  The squared row norms of this deviation
will control the averaged increments.

Identify the current unresolved coordinates with an $l$-element set.
Put
\[
 D=\Xi^TW\Xi-\Pi_l,\quad
 Y=\operatorname{Tr}D,\quad X=\operatorname{Tr}D^2,\quad
 T=\sum_xD_{xx}^2,\quad Z_D=\|D\|.
\]
The analogous reference quantities are denoted with hats.  They are
real symmetric and annihilate $\mathbf1$.  From now on put
$E=W-\widehat P$, so that $D-\widehat D=\Xi^TE\Xi$.
Thus $D$ measures the compressed resolvent's deviation from the identity,
whereas $E$ measures its error against the continued reference.  Small
entries of $E$, together with the reference bounds, keep the pair updates
well defined; the row energies of $D$ will provide the stronger averaged
bounds needed as $l$ decreases.

We track global--global, global--contrast, and contrast--contrast entries
of $E$, using every fixed contrast
$e_{\operatorname{vtx}x}-e_{\operatorname{vtx}y}$ whose two stubs are
still unresolved.  A contrast label is deleted when either of its
endpoints is selected.  These fixed output vectors require no
recentering.  Averaging contrasts recovers centered-stub entries, with
at most a fixed factor in the bounds.

For a selected pair set $I=[e_u,e_v]$, let $F$ be the two-coordinate
swap and let $J_2$ be the all-ones matrix.  We reuse
\[
 B=I_2+\frac{J_2}{l-2},\quad H_0=-\sigma hFB,\quad
 A=I^TDI,\quad K=I^TD^2I,\quad J=I^TD^3I,\quad
 H=(H_0^{-1}+A)^{-1}.
\]
Here $I_2$ is the $2$-by-$2$ identity; the use of $A$ in this
subsection refers only to this block.  The projection onto surviving
contrasts, in the old coordinate space, is
$Q=\Pi_l-\Pi_lIBI^T\Pi_l$.  Since $Q\Pi_lI=0$, the exact
Woodbury formula gives
\begin{equation}
 D'=Q(D-DIHI^TD)Q.
 \label{eq:original-44}
\end{equation}
It gives the same recurrence with hats.  A global row
$w_i=\Xi^TWe_i$ updates to
$Q(w_i-DIHI^Tw_i)$.
Let
$r_i=w_i-c_i\Pi_le_{s(i)}$ while the short endpoint $s(i)$ from
\eqref{eq:original-41} is active, and let $r_i=w_i$ after its removal
or if it did not exist.

The bootstrap on entry errors and the deterministic reference bounds
will give throughout the two phases
\begin{equation}
 \max_{x,y}|D_{xy}|+\max_{i,x}|(r_i)_x|\le\epsilon,
 \qquad \max_{x,y}|\widehat D_{xy}|\le\epsilon,
 \qquad \epsilon=n^{-r_P/2}.
 \label{eq:original-45}
\end{equation}
Changing the displayed right side by a fixed factor makes no
difference, and we do so where contrasts are used.
In particular all two-coordinate Neumann series converge absolutely
with a uniform geometric bound.  The usual inertia argument shows
that every candidate pair preserves positive precision: the ideal
middle matrix $H_0^{-1}$ has one positive and one negative eigenvalue
bounded away from zero, and the perturbation $A$ is $o(1)$.

We give the averaging estimates used below in some detail, to specify
both the cancellations and the absolute bounds.

\begin{lemma}[Coefficient and averaging bounds]
\label{lem:boundary-pair-averages}
Suppose that $D\mathbf1=0$,
$\max_x(D^2)_{xx}\le C_0$ and
$\max_{x,y}|D_{xy}|\le\epsilon=o(1)$.
Then $\|H\|\le C$, and
\[
 |H_{uu}|\le C(l^{-1}+|D_{vv}|+|D_{uv}|).
\]
The conditional absolute mean of $(H-H_0)_{uu}$ with $u$ fixed is
at most $C(1+\sqrt T)/\sqrt l$, and the conditional second moment
of $H_{uu}$ is at most $C(1+T)/l$.
The kernel $|(H-H_0)_{uv}|$ has every row and column sum at most
$C(1+\sqrt T)\sqrt l$.  After subtracting the single-$A$ term
$-H_0AH_0$, these absolute bounds gain a factor $C\epsilon$.
All these statements remain valid if the possible partner set omits
a fixed number of coordinates.

If $x,y$ are arbitrary scalar arrays and $\mathcal I$ is the allowed
endpoint set, $m=|\mathcal I|$, then the pure-swap single-$A$
contraction has the exact formula
\[
 \mathbb E_{\mathcal I}
     [(I^Tx)^TF A F(I^Ty)]
 =\frac{2}{m(m-1)}\left[
   Y_{\mathcal I}\langle x_{\mathcal I},y_{\mathcal I}\rangle
  +x_{\mathcal I}^TD_{\mathcal I\mathcal I}y_{\mathcal I}
  -2\sum_{u\in\mathcal I}D_{uu}x_uy_u\right],
\]
where $Y_{\mathcal I}=\sum_{u\in\mathcal I}D_{uu}$.
For full pairs the following additional identities hold:
\begin{align*}
 \mathbb E\operatorname{Tr}(F A F K)
 &=\frac{2}{l(l-1)}
    \left[YX+\operatorname{Tr}D^3
                 -2\sum_uD_{uu}(D^2)_{uu}\right],\\
 \mathbb E\operatorname{Tr}(F K F K)
 &=\frac{2}{l(l-1)}
    \left[X^2+\operatorname{Tr}D^4
                 -2\sum_u(D^2)_{uu}^2\right],\\
 \mathbb E[(I^Tr)^TF K F(I^Tr)]
 &=\frac{2}{l(l-1)}
    \left[gX+r^TD^2r-2\sum_ur_u^2(D^2)_{uu}\right],
 \qquad g=\|r\|^2.
\end{align*}
\end{lemma}

\begin{proof}
Expand
$H=H_0-H_0AH_0+H_0AH_0AH_0-\cdots$.
Separate $H_0=-\sigma hF+O(l^{-1})$ in each word.
For a diagonal coefficient starting at $u$, a word of pure swaps
with at least one $A$ must touch either $D_{vv}$ or $D_{vu}$ at its
first $A$ position.  Every remaining $A$ factor has norm
$O(\epsilon)$.  Words whose first swap is not pure contribute
$O(l^{-1})$, with the remaining series uniformly bounded.
This proves the pointwise diagonal estimate and, without the
zero-$A$ term, the corresponding estimate for $H-H_0$.
The same first-position argument gives a row bound for the
off-diagonal coefficient; symmetry gives the column bound.
Use
\[
 \sum_v|D_{uv}|\le \sqrt{lC_0},\quad
 \sum_v|D_{vv}|\le\sqrt{lT},\quad
 \sum_v D_{uv}^2\le C_0
\]
to obtain the asserted averages and sums.  A word with at least two
$A$ factors has a further factor $C\epsilon$.  Omitting a bounded
number of coordinates changes these absolute averages by at most a
fixed density factor for large $l$.

For the signed formula expand the two-coordinate expression as
\[
 x_uy_uD_{vv}+x_vy_vD_{uu}
              +(x_uy_v+x_vy_u)D_{uv}
\]
and sum over $u\ne v$ in $\mathcal I$.  The three remaining
identities follow by exactly the same expansion, respectively with
the second matrix $D^2$, with both matrices $D^2$, and with the
rank-one matrix $rr^T$.  No independence of entries is used.
\end{proof}

For later use, the preceding kernel bounds imply the following
bilinear rule.  Under a uniform pair, the absolute contribution of
$H-H_0$ against two outer arrays $x,y$ is at most
\[
 \frac{C(1+\sqrt T)}{l^{3/2}}\|x\|\,\|y\|.
\]
For a same-end term this is the conditional diagonal estimate and
Cauchy--Schwarz; for an opposite-end term it is the row-and-column
sum bound and the elementary Schur estimate.  The remainder after
the single-$A$ term gains $C\epsilon$.  These statements also apply
to restricted arrays after finitely many endpoint omissions.

\par\smallskip\noindent
\textit{The initial energy gap.}
The stop estimates give bounded contrast-row and global-tail energies.
The trace variables transfer directly: $Y=L_1$ at the stop.  For the
diagonal-square statistic $T$, write each $D_{xx}$ as the cleanup
error $(W-P_{<L})_{\xi_x,\xi_x}$ plus the path-reference deviation
from $(\Pi_{l_0})_{xx}$.  The cleanup bound on $T_S$ and the
return-square bound in Lemma~\ref{lemma:path-majorants} therefore
give $T\le n^\eta$, after choosing smaller incoming losses.
The same absolute path majorants give $\widehat T\le n^\eta$.
More importantly, the averaged local identity gives a gap below one in
the normalized boundary energy:
\begin{equation}
 \frac{h^2X}{l_0-1}
 \le1-c_dn^{-\theta}+O(n^{-c_6}),\qquad c_6>\theta.
 \label{eq:original-46}
\end{equation}
Indeed, expand
$X=\operatorname{Tr}(\Xi^TW\Xi)^2-2F_1+(l_0-1)$ and use
\eqref{eq:original-11} to obtain
\[
 \frac{h^2X}{l_0-1}
 =1-\frac{\alpha\mathcal T}{l_0-1}
    -\frac{2h^2L_1}{l_0-1}
    +O(n^{-c_2}+l_0^{-1}).
\]
The cleanup lower bound gives
$\alpha\mathcal T/l_0\ge c_dn^{-\theta}$, since
$l_0/(n\alpha_*)\asymp n^{-.35a}$.
The trace error and the displayed remainder have fixed power margins;
choosing $\theta$ below those margins proves \eqref{eq:original-46}.
The remaining initial bounds needed below are
\[
 T+\widehat T\le n^\eta,\quad
 |Y|\le n^{.09a}/\alpha_*,\quad
 |\widehat Y|\le n^\eta q^{L/2},\quad
 Z_D\le Cn^{v_*}/\alpha_*.
\]

\subsection{Phase A: from averaged energy to small rows}
\label{subsec:boundary-phase-a}

The gap \eqref{eq:original-46} is initially an averaged statement.  We
first make it grow as the boundary shrinks, then use it to contract every
contrast row and every global tail.  Entry estimates are maintained
throughout so that all candidate precisions remain positive.

Choose the fixed exponents $g=r>0$ sufficiently small in the order
specified below, put $\delta\asymp n^{-g}$ and
$l_2\asymp l_0n^{-g-r}$, and set
\[
 \bar u_l=\frac{l}{l+\delta l_0},\quad d_l=1-\bar u_l,
 \qquad \omega_{l_0}=1,
 \quad \omega_{l-2}=\omega_l(1-d_l/l).
\]
Taking logarithms, and summing the $O(l^{-2})$ one-step error, gives
$\omega_l\asymp(\bar u_l/\bar u_{l_0})^{1/2}$.
In particular $\omega_{l_2}\asymp n^{-r/2}$.
We maintain the following bounds simultaneously:
\begin{equation}
 \begin{gathered}
 Z_D\le\sqrt l\,n^{-.1a},\qquad
 |Y|+|\widehat Y|\le2\sqrt{l_0}\,n^{-.04a},
 \qquad T+\widehat T\le n^\eta,\\
 u_l:=\frac{h^2X}{l-1}\le\bar u_l+o(n^{-g}),\\
 \|D(e_x-e_y)\|^2\le C_S\omega_l,
 \qquad \|r_i\|^2\le C_V\omega_l,\\
 \max|E_{ij}|\le C_EB_A,\qquad
 E=W-\widehat P,\qquad B_A=\frac{n^{.38a}}{\sqrt{l_2}}.
 \end{gathered}
 \label{eq:original-47}
\end{equation}
The last maximum is over the three types of tracked output pairs.
The constants in the row and entry bounds can take finitely many
values, as specified at the exceptional events below.  Initially
they are chosen above the bounded initial energies.  The row and entry
account thresholds leave fixed slack inside these caps, to accommodate
the power-small excesses obtained by extraction.  The contrast
bound gives an individual row bound because
$De_x=l^{-1}\sum_yD(e_x-e_y)$ and Jensen's inequality applies.

\subsubsection*{Operator norm control}
The first estimate in \eqref{eq:original-47} keeps the operator norm
small relative to $\sqrt l$.  A fixed spectral moment suffices; the much
higher moments used later for individual entries are not needed here.
Fix the even integer $k=8$.
Orthogonal projection contracts every Schatten norm.  In the expansion
of $\operatorname{Tr}(D-DIHI^TD)^k$, the words with exactly one update
factor sum to $-k\operatorname{Tr}[H(I^TD^{k+1}I)]$.
The $H_0$ contribution has zero mean by centering.  For the other part,
\[
 |(D^{k+1})_{xy}|
 \le Z_D\sqrt{(|D|^k)_{xx}(|D|^k)_{yy}},
\]
so Lemma~\ref{lem:boundary-pair-averages} gives mean absolute value
at most
$C_k(1+\sqrt T)(Z_D/\sqrt l)\operatorname{Tr}|D|^k/l$.

A word with $p\ge2$ update factors has, after a cyclic permutation
under the trace, gaps $a_1,\ldots,a_p\ge0$ with
$\sum_ja_j=k-p$.  Its absolute trace is at most
$C_k\prod_jm_{2+a_j}$, where
$m_s=\operatorname{Tr}(I^T|D|^sI)$.
Indeed $-|D|^s\preceq D^s\preceq |D|^s$, so the norm of
$I^TD^sI$ is at most $m_s$.  The product acts in dimension two.  Interpolation of the
positive spectral measure induced by the two selected coordinates
gives
$m_{2+a}\le m_2^{1-a/(k-2)}m_k^{a/(k-2)}$.
Here $m_2\le2C_0$ and
$\sum_ja_j/(k-2)=(k-p)/(k-2)\le1$.  Consequently the product is
at most $C_k(1+m_k)$.
Averaging $m_k$ gives $2\operatorname{Tr}|D|^k/l$.  Under the
bootstrap, $(1+\sqrt T)Z_D/\sqrt l=o(1)$, and hence
\[
 \mathbb E\operatorname{Tr}|D'|^k
 \le(1+C_k/l)\operatorname{Tr}|D|^k+C_k.
\]
Use the account with normalization
$l_0(n^{v_*}/\alpha_*)^{k-2}$.
Its initial value is bounded by the initial row and norm estimates;
the additive terms over all steps are bounded after this normalization.
The discount loss is at most $n^{C_k'(g+r)}$.
Before this loss and the change from $l_0$ to $l$, its $k$th-root
bound relative to $\sqrt{l_0}$ is
\[
 \left(\frac{n^{v_*}}{\alpha_*\sqrt{l_0}}\right)^{1-2/k}
 =n^{-.45a(1-2/k)+o(1)}.
\]
Taking $g+r$ sufficiently small proves the norm bound in
\eqref{eq:original-47} with a strict power margin.

\subsubsection*{Trace and diagonal estimates}
We next control the traces appearing in signed averages.  The diagonal
statistic requires a separate argument: its growth must have an
arbitrarily small power loss.

Expanding \eqref{eq:original-44}, using $Q^2=Q$ and cyclicity of
the trace, gives the following exact trace identities:
\begin{equation}
 \begin{split}
 \Delta Y={}&-\operatorname{Tr}BA
                  -\operatorname{Tr}H(K-ABA),\\
 \Delta X={}&-2\operatorname{Tr}BK+\operatorname{Tr}BABA\\
 &-2\operatorname{Tr}H(J-ABK-KBA+ABABA)\\
 &+\operatorname{Tr}H(K-ABA)H(K-ABA).
 \end{split}
 \label{eq:original-48}
\end{equation}
The factors follow from
$Q=\Pi_l-\Pi_lIBI^T\Pi_l$ and $D\Pi_l=D$; in particular
$I^TDQD I=K-ABA$ and
$I^TDQDQDI=J-ABK-KBA+ABABA$.

For $Y$, centering gives
$\mathbb E\operatorname{Tr}BA=2Y/(l-1)$ and
$\mathbb E\operatorname{Tr}H_0K=0$.
The single-$A$ term contributes
$2h^2YX/[l(l-1)]+O((1+Z_D)/l)$ by the first full-pair identity
of Lemma~\ref{lem:boundary-pair-averages}.
The remaining terms have mean absolute value at most $C(T+1)/l$:
the pair $K$ is bounded, and
$\mathbb E\|A\|^2\le C(T+1)/l$.
Also
\[
 \mathbb E(\Delta Y)^2\le C(Z_D^2+T+1)/l.
\]
For the off-diagonal $HK$ term use
$\operatorname{Tr}D^4\le Z_D^2X\le CZ_D^2l$; for its diagonal
term use the conditional second-moment coefficient bound; the
remaining terms use $\mathbb E\|A\|^2$ and bounded $K,H$.
Thus the drift is a coefficient bounded in magnitude by $C/l$
times $Y$, plus absolute error $C(Z_D+T+1)/l$.
The hatted argument is identical and uses $\|\widehat D\|\le C/R_0$.
Square accounts at scale $\sqrt{l_0}n^{-.06a}$ have initial bounded
values.  Indeed $\alpha_*\sqrt{l_0}\asymp n^{.575a}$ and
$q^{L/2}/\sqrt{l_0}\asymp n^{-.29a}$.
Their mean-error and variance budgets have a strict power margin at
this scale.  A discount loss $n^{C(g+r)}$ is absorbed before
extracting the stated trace caps.

We control $T$ by retaining the loss from deleted diagonals.  For a surviving coordinate $x$, put
$f_x=\Pi_le_x$ and $a_x=I^TDe_x$.
Its new centered vector is
$f_x+(f_u+f_v)/(l-2)$.  The old-vector diagonal increment is
$-a_x^THa_x$; the centering correction has pointwise magnitude
$C\epsilon/l$ and conditional absolute mean $C/l^{3/2}$.
For the latter assertion, linear correction terms use
$\sum_u|D_{xu}|\le C\sqrt l$; update terms use bounded pair $H$
and the same row bound, and terms with two centering shifts carry
$l^{-2}$.  Conditional on survival of $x$, the $H_0$ core mean is
$O(l^{-2})$ by the centered-pair formula.  The single-$A$ core mean
is bounded by
$C(|Y|+Z_D+1)/l^2$ by the restricted single-$A$ identity, whereas
the higher remainder is bounded by
$C\epsilon(1+\sqrt T)/l^{3/2}$ by the bilinear rule.
Under \eqref{eq:original-47},
$(|Y|+Z_D)/\sqrt l$ is power-small.

For a fixed $x$, the second moment of the core increment is at most
$C(1+\epsilon^2T)/l^2$.  Opposite-end terms use the product of the
two row squared sums, divided by $l^2$.  Same-end terms use the
conditional bound on $H_{uu}^2$ and
$\sum_uD_{xu}^4\le\epsilon^2\sum_uD_{xu}^2\le C\epsilon^2$.
The centering squares contribute at most $C\epsilon^2/l$ when
summed over $x$.  Selection deletes diagonals with expected square
mass exactly $2T/l$.  Finally, multiplying each conditional mean
bound by $2|D_{xx}|$ and using
$\sum_x|D_{xx}|\le\sqrt{lT}$ shows that the signed first-order
cost is at most $C(\sqrt T+\epsilon T)/l$.
Young's inequality absorbs this and the variance terms into a fixed
portion of the deletion term, leaving
\begin{equation}
 \mathbb ET'\le T+C/l.
 \label{eq:original-49}
\end{equation}
The proof with hats is the same.  Normalize these accounts by a
slightly larger arbitrarily small power of $n$ than their initial
losses; since $\sum l^{-1}=O(\log n)$, the total additive budget is
bounded after this normalization.  Extraction with another small
power cushion gives $T+\widehat T\le n^\eta$.

\subsubsection*{Growth of the averaged gap}
We now compare the normalized energy $u_l$ with $\bar u_l$.
For $X$, the ideal averages in \eqref{eq:original-48} are
$-4X/(l-1)$, zero from $-2H_0J$, and
$2h^2X^2/[l(l-1)]$ from the quadratic term, with error
$C(1+Z_D^2)X/l^2$.  The full-pair identities above give these
coefficients directly.  Furthermore
$|J_{xy}|\le Z_D\sqrt{(D^2)_{xx}(D^2)_{yy}}$; the bilinear
coefficient rule bounds the $\Delta H J$ mean by
$C(1+\sqrt T)(Z_D/\sqrt l)X/l$.
All remaining errors are bounded by
$C(\epsilon X/l+(T+1)/l)$: terms with one small $A$ and pair $K$
use $\|A\|\le C\epsilon$; terms with two $A$ factors use
$\mathbb E\|A\|^2\le C(T+1)/l$.
The second-moment bound is
\[
 \mathbb E(\Delta X)^2
 \le C\{1+Z_D^4/l+(1+T)Z_D^2/l\}.
\]
Only the $HJ$ term needs the last two quantities.  For its
off-diagonal part use
$\operatorname{Tr}D^6\le Z_D^4X\le CZ_D^4l$; for the diagonal
part use $|J_{uu}|\le C Z_D$ and the conditional bound for
$H_{uu}^2$.  The other terms have bounded second moment because
their pair factors are bounded.
Including the changed denominator in $u_l=h^2X/(l-1)$ gives
\begin{equation}
 \mathbb E\Delta u_l
   \le-\frac2l u_l(1-u_l)+\frac{n^{-c_7}}l,
 \qquad
 \sum_{l_2\le l\le l_0}\mathbb E(\Delta u_l)^2
   \le n^{-c_8}
 \label{eq:original-50}
\end{equation}
for fixed positive $c_7,c_8$ after sufficiently small loss choices.
For example the $Z_D^4/l$ part of the last bound contributes at
most $Cn^{-.4a}\sum l^{-1}$.  The error terms in the first bound
are powers because $Z_D/\sqrt l\le n^{-.1a}$ and
$\epsilon=n^{-r_P/2}$.

The comparator satisfies
$\bar u_{l-2}=\bar u_l-2\bar u_l(1-\bar u_l)/l+O(l^{-2})$.
The drift of a positive excess $u_l-\bar u_l$ is at most
$C/l$ times that excess plus $Cn^{-c_7}/l$.
Apply the square positive-excess account at scale
$n^{-\zeta_X}$, with
\[
 g+C(g+r)<\zeta_X<\tfrac14\min(c_7,c_8).
\]
Such a choice exists by first making $g,r$ smaller.  Its discount
loss is at most $n^{C'(g+r)+o(1)}$, and its initial positive excess
vanishes by \eqref{eq:original-46}, since $\theta\ll g$.
The inequality
$(x+y)_+^2\le x_+^2+2x_+y+y^2$ and
\eqref{eq:original-50} therefore give
$u_l\le\bar u_l+o(n^{-g})$, with a power margin in the little-oh
term.

\subsubsection*{Contraction of individual rows}
The averaged gap supplies a negative homogeneous term in each surviving
row's energy drift.
Contrast rows have a homogeneous update.  A global tail has one extra
centering term while its short endpoint survives, and one exceptional
transition when that endpoint is removed.

Let $f=e_x-e_y$ be a surviving contrast, and put
$r_f=Df$, $g_f=\|r_f\|^2$, $a_f=I^Tr_f$ and
$b_f=I^TDr_f$.  Since $Qf=f$ on its survival event, expansion of
\eqref{eq:original-44} gives
\begin{equation}
 \Delta g_f=-a_f^TBa_f-2a_f^THb_f
       +2a_f^THABa_f+a_f^TH(K-ABA)Ha_f.
 \label{eq:original-51}
\end{equation}
More generally the same identity describes the energy of a centered
row updated homogeneously by $r'=Q(r-DIHI^Tr)$.
Write $g=\|r\|^2$, $a=I^Tr$ and $b=I^TDr$.
Suppose its entries are bounded by $C\epsilon$ and $g\le C$.
Conditioning the pair to omit a fixed number of coordinates where
this same entry bound holds gives the following estimates.

First, deletion has upper mean $-2g/l+C\epsilon^2/l$.
The full $a^TH_0b$ mean vanishes; on the restricted set the
centered-pair formula bounds it by
$C/l^2+CZ_Dg/l^3$.  Here each omitted $r_s$ is $O(\epsilon)$,
each omitted $(Dr)_s$ is bounded by
$\sqrt{(D^2)_{ss}g}$, and the restricted dot product is at most
$Z_Dg$ plus a bounded omission correction.
The $\Delta H$ contribution is at most
$C(1+\sqrt T)Z_Dg/l^{3/2}$ by the bilinear rule applied to
$r$ and $Dr$.
The term $a^THABa$ has upper absolute mean $C\epsilon g/l$.
Finally the ideal quadratic term is nonnegative, so restricting the
pair inflates its full mean by at most $1+O(l^{-1})$.
The last identity of Lemma~\ref{lem:boundary-pair-averages} bounds
that full mean by
\[
 \frac{2h^2gX}{l(l-1)}
       +\frac{C(1+Z_D^2)g}{l^2}.
\]
Replacing $H$ by $H_0$, or $K-ABA$ by $K$, costs
$C\epsilon g/l$ using bounded pair $K$ and
$\|H-H_0\|\le C\epsilon$.

The leading deletion and quadratic terms therefore combine as
\[
 -\frac{2g}{l}+\frac{2h^2gX}{l(l-1)}
 =-\frac{2g}{l}(1-u_l).
\]
This is where the averaged gap becomes a contraction for each row.
Under the Phase A caps,
$1-u_l\ge(1-o(1))d_l$.  The error coefficients just listed
are $o(d_l)$, since $d_l\ge cn^{-g}$ and $g$ is chosen below
all their fixed exponent margins.  We obtain
\begin{equation}
 \mathbb E_{\rm surv}\Delta g
 \le-\frac{1.5d_lg}{l}+\frac{n^{-c_9}}l
 \label{eq:original-52}
\end{equation}
for a fixed $c_9>0$.
For the second moment of the homogeneous increment,
\begin{equation}
 \mathbb E_{\rm surv}(\Delta g)^2
 \le C\left(
       \frac{\epsilon^2g}{l}
       +\frac{Z_D^2g^2}{l^2}
       +\frac{n^\eta\epsilon^2Z_D^2g}{l^2}
       \right).
 \label{eq:original-53}
\end{equation}
To verify this bound, first consider the terms in
\eqref{eq:original-51} that do not involve $b$.  Each is bounded
by $C\|a\|^2$, so its square averages to
$C\epsilon^2g/l$.  The opposite-end part of $a^THb$ uses
$\|r\|^2\|Dr\|^2/l^2\le Z_D^2g^2/l^2$.
For its same-end part, use the conditional bound
$\mathbb EH_{uu}^2\le C(1+T)/l$, the pointwise bound
$|r_u|^2\le C\epsilon^2$, and
$\sum_u(Dr)_u^2\le Z_D^2g$.
The finite sum of these terms changes constants only.
The same identity also gives the pointwise bound $|\Delta g|\le
C\epsilon$, since $|b_s|\le\sqrt{(D^2)_{ss}g}\le C$.

For a global tail with an active short endpoint, an ordinary event
omits that endpoint.  Its exact update is
\[
 r_i'=Q\{r_i-DIH(I^Tr_i-k_i\mathbf1_2)\},\qquad k_i=c_i/l.
\]
If its short endpoint is absent, put $k_i=0$.
To estimate the additional energy terms, write the additional vector
as $k_iQDIH\mathbf1_2$.
Its cross term with $Qr_i$ involves
$I^TDQr_i=b_i-ABa_i$.
Averaging its absolute value gives at most
$C Z_D\sqrt g/\sqrt l+C\sqrt g/\sqrt l$ before multiplication
by $|k_i|\le l^{-1}$.
Its cross term with $QDIHa_i$ is bounded in mean by
$C\sqrt g/l^{3/2}$, using bounded pair $K$ and
$\mathbb E\|a_i\|\le C\sqrt{g/l}$.
The square of the additional vector is $O(l^{-2})$.
Thus the additional mean budget is
\[
 C Z_D\sqrt g/l^{3/2}+C/l^{3/2}+C/l^2.
\]
It is power-small beyond $l^{-1}$ under the Phase A norm bound.
Its pointwise effect on the energy is $O(l^{-1})$, so it adds
$C/l^2$ to the second-moment budget.  Consequently
\eqref{eq:original-52}--\eqref{eq:original-53} remain sufficient
for ordinary global-tail steps, after decreasing $c_9$ and adding
this summable term.

If $s(i)$ is selected, it disappears permanently and
$Q\Pi_le_{s(i)}=0$.  The exact update instead gives the pathwise
bound
\begin{equation}
 \|r_i'\|\le\|r_i\|+C\sqrt{\max_s(D^2)_{ss}}.
 \label{eq:original-54}
\end{equation}
This is a single exceptional event for that row.

Apply upper positive-excess accounts to $g/\omega_l$ above a fixed
constant.  Combining \eqref{eq:original-52} with
$\omega_{l-2}/\omega_l=1-d_l/l$ leaves negative homogeneous drift;
the remaining upper mean is at most $n^{-c_{10}}/l$.
The sums of the second moments after division by $\omega_l^2$
are power-small.  Explicitly, $\omega_l\ge cn^{-r/2}$ gives
\[
 \sum\frac{\epsilon^2g}{l\omega_l^2}
       \le C\epsilon^2n^{r/2}\log n,
 \qquad
 \sum\frac{Z_D^2g^2}{l^2\omega_l^2}
       \le Cn^{-.2a}\log n;
\]
the third term of \eqref{eq:original-53} is bounded by the product
of these small factors up to $n^\eta$, and the additional $l^{-2}$
terms sum to $O(n^r/l_2)$.
The change of normalization contributes only $O(l^{-2})$ to this
normalized second budget.
Normalized pointwise increments are power-small because
$\epsilon n^{r/2}=o(1)$.
The positive-excess account lemma, with a sufficiently small scale
exponent and then a sufficiently large even order, controls all the
polynomially many contrast and global labels simultaneously.
Contrast labels die at endpoint selection.  At the exceptional
global event \eqref{eq:original-54}, the contrast bound gives
$\max_s(D^2)_{ss}\le C_S\omega_l$; thus a larger fixed
post-event threshold makes its positive excess zero.
There is at most one such threshold increase per global row.
This proves the row-energy line of \eqref{eq:original-47}.

\subsubsection*{Entry comparison in Phase A}
It remains to control $E=W-\widehat P$.  For an entry whose short
coordinates are not selected, both actual and reference outer rows are
small at the chosen endpoints.  This gives concentration.  The finitely
many selections of short coordinates are handled separately below.

For any two fixed surviving output vectors, the numerical update is
\begin{equation}
 \Delta E_{ij}
   =-(I^Tw_i)^TH(I^Tw_j)
         +(I^Tp_i)^T\widehat H(I^Tp_j).
 \label{eq:original-55}
\end{equation}
For a contrast $f$, the rows here are
$w_i=f+Df$ and $p_i=f+\widehat Df$.
Call an event ordinary for this entry when its contrast labels
survive and none of its active global short endpoints is selected.
This omits at most a fixed number of coordinates from the uniform
pair.  On the allowed endpoint set, all actual and reference outer
entries are $O(\epsilon+l^{-1})$, their full row norms are bounded,
and their full sums are zero.  The difference outers and the block
$D-\widehat D$ have entry bounds $CB_A$ by the induction hypothesis.

The conditional absolute mean budget is
\begin{equation}
 |\mathbb E_{\rm ord}\Delta E_{ij}|
 \le C\left\{
   \frac{1+|Y|+|\widehat Y|+Z_D}{l^2}
       +\frac{B_A\epsilon}{l}+\frac{n^\eta}{l^{3/2}}
       \right\}.
 \label{eq:original-56}
\end{equation}
To prove this, expand both coefficients in \eqref{eq:original-55}
around $H_0$.  The zero-$A$ terms cancel in full-pair expectation.
For the finitely many omitted coordinates, the centered-pair formula
and bounded row entries give $C/l^2$.
For the single-$A$ terms use the signed restricted identity in
Lemma~\ref{lem:boundary-pair-averages}.  Its diagonal factor is
$|Y|+O(1)$, its cross factor is at most $Z_D$, and the outer
dot products are bounded.  This gives the first fraction, and the
hatted version uses $\|\widehat D\|\le C$.
Replacing a pure swap by its $O(l^{-1})$ part gives smaller terms,
since an outer same-end product has average $O(l^{-1})$ and an
opposite-end product is bounded by the row norms.

For words with at least two $A$ factors, telescope the difference
of the two products.  A changed endpoint factor has size $CB_A$,
and another endpoint factor contributes $C\epsilon$; averaging
the outer product costs $C/l$.  This gives $CB_A\epsilon/l$.
If all endpoint factors are hatted, the difference is in an outer
row.  Its norm is bounded, and the hatted remainder kernel after
the single-$A$ term is at most
$C\epsilon(1+\sqrt{\widehat T})/l^{3/2}$ against bounded row
norms.  This is at most the last fraction of
\eqref{eq:original-56}.  Summation over longer words is justified
by their common geometric factor $C\epsilon$; the number of
possible locations of a specified changed factor is absorbed into
that summation.

The conditional second moment is at most
$C\{1+n^\eta(\epsilon+l^{-1})^2\}/l^2$.
For opposite-end terms this is the product of two bounded row
squared sums divided by $l^2$.
For same-end terms use the allowed outer supremum and the
conditional bound on $H_{uu}^2$ from the coefficient lemma; this
gives the additional displayed factor.  Bounding the two terms of
\eqref{eq:original-55} separately is sufficient for this estimate.
For the pointwise difference, use
\[
 H-\widehat H=-H(A-\widehat A)\widehat H
\]
and telescope its outer factors.  It gives
$|\Delta E_{ij}|\le CB_A(\epsilon+l^{-1})$ on ordinary events.

The mean budget in \eqref{eq:original-56}, divided by $B_A$, and
the second budget divided by $B_A^2$, have power-small sums.
For example, the trace contribution to the first sum is bounded by
\[
 \frac{C\sqrt{l_0}n^{-.04a}}{l_2B_A}
   =Cn^{-.42a+(g+r)/2},
\]
the last term contributes at most $Cn^{\eta-.38a}$, and the
second-moment sum is $O(n^{-.76a+o(1)})$.
The norm term and the finite-size terms have still smaller bounds.
The middle mean term gives $C\epsilon\log n$.
These explicit comparisons justify the required choice of the
positive-excess scale exponent independently of its later moment
order.

A surviving exceptional event selects a global short endpoint.  Put $e_i=w_i-p_i$.
An exact telescoping form of \eqref{eq:original-55} is
\[
 \Delta E_{ij}
 =-(I^Te_i)^TH(I^Tw_j)
  -(I^Tp_i)^TH(I^Te_j)
  +(I^Tp_i)^TH(A-\widehat A)\widehat H(I^Tp_j).
\]
For a global--contrast entry, the contrast endpoint coordinates
are omitted on every survival event, so its actual and reference
outer entries are small.  The only unsuppressed input error in
this display is a contrast--contrast error; the coefficient on an
input global--contrast error is $O(\epsilon+l^{-1})$.
For a global--global entry, all input errors on the right are of
global--contrast or contrast--contrast type.  Thus choose the
constant thresholds in the following finite order: first the
contrast--contrast threshold, then pre-event and post-event
global--contrast thresholds, and then global--global thresholds
indexed by the number, zero, one or two, of short-endpoint removals.
Increase each next threshold by a fixed amount exceeding the jump
bound supplied by already chosen lower-type thresholds.  For large
$n$ the small feedback factor in the global--contrast case is
absorbed by this cushion.  Each exceptional event then resets the
relevant positive and negative excesses to zero.

Initial entries are $o(B_A)$ by the cleanup bound, the taper-change
bound and the fixed contrast combinations.
Apply the positive-excess account lemma on ordinary events to both
signs of every entry.  Keep the family denominators fixed and drop
dead contrast labels.  The moment order is fixed large enough to
overcome the polynomial number of labels; the ordinary discount
budgets are summably power-small.  The exceptional threshold rule
is pathwise and costs no discount.  This proves the entry line of
\eqref{eq:original-47}, and hence also
\eqref{eq:original-45}.

\subsubsection*{Simultaneous choice}
The norm, trace, averaged-energy, row and entry estimates have all been
derived under the current bounds \eqref{eq:original-47}.  We now use them
together to obtain those bounds at the next chosen state.

All the preceding accounts are placed in the same deterministic
sum.  Its initial value is $n^{o(1)}$: the scalar accounts are bounded
after normalization, each polynomial family of individual excess accounts
has its fixed reciprocal-family weight and initially vanishing excess,
and there are only $O(\log n)$ retained spectral parameters.
The only permitted polynomial discount losses in Phase A
are the fixed spectral-moment, trace-square and scalar-comparator
losses $n^{C(g+r)}$.  Their constants do not depend on the large
moment orders later used for maximum extraction.  First fix the
initial row comparator constants and the spectral order $8$;
then choose $g=r$ sufficiently small for those polynomial losses,
the initial gap \eqref{eq:original-46}, and the fixed margins
$a,r_P$.  Choose the positive-excess scale exponents next and their
large even orders last.  Arbitrarily small structural and diagonal
losses are made smaller as necessary.  This gives a joint induction
for \eqref{eq:original-47}.  In particular, for a fixed
$0<r_g<r/2$ and all sufficiently large $n$, every surviving
contrast squared energy and every global-tail squared energy at
$l_2$ is at most $n^{-r_g}$.

\subsection{Phase B: entry concentration on the shrinking boundary}
\label{subsec:boundary-phase-b}

At $l_2$ every contrast row and global tail has power-small energy.  We
allow slightly larger energy bounds and use this margin to control entry
errors down to $l_3$.  Bounds on $Y$ and $\widehat Y$ are no longer
needed; we retain the diagonal-square and entry estimates.  Summing the
Phase A entry second-moment bound to $l_3$ would give only $O(l_3^{-1})$,
which is too large for the desired scale $B_B^2$.  We will instead use
the small row energies when estimating the difference of the actual and
reference updates.

Fix $0<r_v<r_s<r_g$, with strict power gaps, and put
$G_S=n^{-r_s}$ and $G_V=n^{-r_v}$.
The new barriers allow fixed multiples of these squared-energy
scales.  Choose $\chi>r_s,r_v$ as in the parameter hierarchy and
let $l_3\asymp n^\chi$.  From $l_2$ down to $l_3$ maintain
\[
 \|D(e_x-e_y)\|^2\le C G_S,\qquad
 \|r_i\|^2\le C G_V,\qquad
 T+\widehat T\le n^\eta,\qquad
 \max|E_{ij}|\le C B_B,
 \quad B_B=\frac{n^{.40a}}{\sqrt{l_2}}.
\]
The constants for global and entry labels again admit the finite
exceptional-event enlargements described above.  All initial values
have power cushions.  In particular \eqref{eq:original-45} remains
available.  The individual row bound implied by contrasts gives
\begin{equation}
 Z_D\le\|D\|_{\rm F}\le C\sqrt{lG_S}.
 \label{eq:original-57}
\end{equation}
The deterministic hatted row bounds are smaller than this scale.

\subsubsection*{Diagonal squares without a trace cap}
We retain \eqref{eq:original-49} in this phase without controlling
$Y$.  For a fixed surviving coordinate the pure $H_0$ core mean
is $O(l^{-2})$ by centering.  Every other coefficient is bounded
absolutely by the bilinear coefficient rule against a row with
squared norm $CG_S$, giving
$C(1+\sqrt T)G_S/l^{3/2}$.
The conditional absolute mean of the centering correction is now
$C\sqrt{G_S}/l^{3/2}+C\epsilon/l^2$.
This follows from $\sum_u|D_{xu}|\le C\sqrt{lG_S}$ in the
linear correction and from the entry bound in the remaining
two-shift terms.

The second moments, summed over surviving diagonals, are bounded
by
\[
 \frac C l\{G_S^2+\epsilon^2G_S(1+T)+\epsilon^2\}.
\]
For opposite-end core terms use
$\sum_x(\sum_uD_{xu}^2)^2\le C lG_S^2$ followed by the
$l^{-2}$ pair average.  Same-end terms use the conditional bound
for $H_{uu}^2$ and
$\sum_{x,u}D_{xu}^4\le C\epsilon^2lG_S$.
The last term covers the centering squares.
The signed first-order cost, after multiplying the preceding
conditional means by $2|D_{xx}|$ and summing, is at most
\[
 \frac C l
    \{G_S(1+\sqrt T)+\sqrt{G_S}\}\sqrt T
       +C\sqrt{lT}/l^2.
\]
Deletion subtracts $2T/l$.  Since $G_S$ and $\epsilon$ tend to
zero, the terms proportional to $T/l$ are absorbed directly;
Young's inequality absorbs the terms proportional to
$\sqrt T/l$, leaving $C/l$.  The hatted proof is identical.
This proves \eqref{eq:original-49} again and hence the new
arbitrarily small-loss barriers for $T,\widehat T$.

\subsubsection*{Preserving the small row energies}
For a surviving contrast apply \eqref{eq:original-51}.
Now $X/l\le CG_S$, and the relative error in the
$\Delta H$ cross term is
$Cn^\eta Z_D/\sqrt l\le Cn^\eta\sqrt{G_S}=o(1)$.
The quadratic finite-size error $Z_D^2/l$ and the small coefficient
errors are also $o(1)$.  The omitted-coordinate error is at most
$C\epsilon^2/l$, as before.  Consequently
\begin{equation}
 \mathbb E_{\rm surv}\Delta g_f
       \le-cg_f/l+C(\epsilon^2/l+l^{-2})
 \label{eq:original-58}
\end{equation}
with a fixed $c>0$.
Using $g_f\le CG_S$ and \eqref{eq:original-57}, the sum of
\eqref{eq:original-53} divided by $G_S^2$ is at most
\[
 C(\epsilon^2/G_S+G_S+n^\eta\epsilon^2)\log n=o(1)
\]
by a power.  The normalized additive drift has sum
$O((\epsilon^2/G_S)\log n+1/(l_3G_S))$, also power-small.
For a pointwise bound,
$|a_f|\le C\epsilon$ and
$|(Dr_f)_u|\le C G_S$; hence the increments in
\eqref{eq:original-51}, divided by $G_S$, are power-small.
The upper positive-excess account at the loosened scale $G_S$
therefore keeps all contrast barriers.

For an ordinary global-tail step the homogeneous calculation is
the same, with $g\le CG_V$.  The extra shift from $k_i=c_i/l$
has mean absolute cross budget
\[
 C\sqrt{G_S}\sqrt g/l+CG_S/l^2+C\sqrt g/l^{3/2}.
\]
Here use $\|Dr\|/\sqrt l\le C\sqrt{G_S}\sqrt g$ and
$(D^2)_{ss}\le CG_S$ in the exact shift calculation from Phase A.
Young's inequality against a portion of the negative homogeneous
drift leaves upper drift
$C(G_S+\epsilon^2)/l+C/l^2$.
Its sum divided by $G_V$ is power-small because
$r_v<r_s$ and $l_3G_V\longrightarrow\infty$.
The homogeneous variance divided by $G_V^2$ has the bound
\[
 C(\epsilon^2/G_V+G_S+n^\eta\epsilon^2G_S/G_V)\log n,
\]
which is power-small.  The shift has pointwise energy effect at
most $C(\sqrt{G_Sg}+G_S)/l$ and contributes a smaller summable
second budget.  Homogeneous normalized increments are small by
$|a|\le C\epsilon$ and
$|b_u|\le C\sqrt{G_SG_V}$.
Thus the same account lemma keeps the global barrier.  The single
exceptional event \eqref{eq:original-54} is absorbed by a fixed
threshold increase; its additional squared scale is $G_S\ll G_V$.

\subsubsection*{Entry concentration down to $l_3$}
For an ordinary event of a tracked entry let $\mathcal I_o$ be its
allowed endpoint set.  The outer arrays in
\eqref{eq:original-55} satisfy
\begin{equation}
 \begin{split}
 \|w_{\mathcal I_o}\|^2+\|p_{\mathcal I_o}\|^2
      &\le C W_g,\qquad W_g=G_V+l^{-1},\\
 \|(w-p)_{\mathcal I_o}\|^2
      &\le C E_g,\qquad E_g=\min(G_V,lB_B^2).
 \end{split}
 \label{eq:original-59}
\end{equation}
For contrasts the identity coordinates have been omitted; for
globals the short coordinate has been omitted and its remaining
centering constant has squared norm $O(l^{-1})$.
The actual tails have squared norm $CG_V$ and the reference tails
are smaller by \eqref{eq:original-43}.
Their difference has no short-coordinate term and has entry
supremum $CB_B$, giving both bounds in the minimum.
Every allowed outer entry is $O(\epsilon+l^{-1})$.
Because the full arrays are centered and only boundedly many
coordinates are omitted, their sums on $\mathcal I_o$ are $O(1)$;
the difference-array sums are $O(B_B)$.

We verify the ordinary conditional mean with these finer norms.
For the ideal $H_0$ terms, the centered-pair formula gives
$CB_B/l^2$.  Its sum-product difference uses the preceding
$O(1)$ and $O(B_B)$ bounds.  Its dot-product correction has a
coefficient $O(l^{-3})$ and a difference bounded by
$C\sqrt l B_B\sqrt{W_g}$, which is smaller.

For single-$A$ pure-swap terms telescope the endpoint block and
the two outer rows.  The endpoint error block has diagonal sum
and operator norm at most $ClB_B$.  The signed restricted
single-$A$ identity and \eqref{eq:original-59} therefore give
$CB_BW_g/l$ for an endpoint-block difference.
For an outer-row difference with a hatted endpoint block,
the diagonal sum of that block is at most
$\sqrt{l\widehat T}$ and the dot difference is at most
$C\sqrt lB_B\sqrt{W_g}$.
After division by the pair count this is
$Cn^\eta B_B\sqrt{W_g}/l$.
The cross contraction uses $\|\widehat D\|\le C$ and gives
the smaller bound $CB_B\sqrt{W_g}/l^{3/2}$.
The diagonal exclusions in the exact identity are smaller by the
same restricted dot bounds.  An $O(l^{-1})$ swap correction is
also smaller: use either the endpoint-error supremum or a hatted
endpoint supremum together with the outer difference, and then
the extra $l^{-1}$ coefficient.

For a word with at least two endpoint factors, a telescoping
endpoint difference contributes
$CB_B\epsilon W_g/l$.
If all endpoint factors are hatted, the remainder coefficient
kernel has bound $Cn^\eta\epsilon/l^{3/2}$; against outer norms
$C\sqrt lB_B$ and $C\sqrt{W_g}$ this gives
$Cn^\eta\epsilon B_B\sqrt{W_g}/l$.
The geometric series over additional endpoint factors converges
uniformly.  Combining these estimates, the sum of absolute
conditional means divided by $B_B$ is bounded by
\begin{equation}
 C\sum_{l\ge l_3}
 \left\{\frac1{l^2}
   +\frac{G_V+l^{-1}+n^\eta(\sqrt{G_V}+l^{-1/2})}{l}\right\}.
 \label{eq:original-60}
\end{equation}
This is at most a constant times
$l_3^{-1}+G_V\log n+n^\eta\sqrt{G_V}\log n
 +n^\eta l_3^{-1/2}$, and is power-small when $\eta$ is chosen
below the fixed gaps.

The conditional second moment in \eqref{eq:original-55} is bounded by
the sum of three terms:
\[
 \frac{C(E_gW_g+B_B^2W_g^2)}{l^2},\qquad
 \frac{Cn^\eta B_B^2W_g}{l^2},\qquad
 \frac{CB_B^2(\epsilon+l^{-1})^2W_g}{l}.
\]
These are respectively: opposite-end terms after telescoping an
outer or the coefficient; same-end terms with an outer error,
using its entry supremum and the conditional square of $H_{uu}$;
and same-end terms with a coefficient error, using
$\|H-\widehat H\|\le CB_B$ and the allowed outer supremum.
There are only finitely many two-coordinate terms, so squaring
their sum changes a fixed constant.
After division by $B_B^2$, the first numerator uses
$E_g\le ClB_B^2$ and its sum is bounded by
$C\sum W_g/l+ C\sum W_g^2/l^2$.
The other two sums are smaller by the displayed factors.  All
are power-small by the same calculation as
\eqref{eq:original-60}.
The pointwise ordinary bound is again
$|\Delta E_{ij}|\le CB_B(\epsilon+l^{-1})$.

Use positive-excess accounts for both signs at the loosened scale
$B_B$.  Their normalized ordinary mean, variance and jump budgets
are power-small.  The same finite exceptional-event threshold
hierarchy as in Phase A applies verbatim to the exact telescoping
identity; surviving contrasts still exclude their identity
coordinates.

To extract all the Phase B barriers jointly, reinitialize the row and
entry accounts at $l_2$ with the loosened scales $G_S,G_V,B_B$, and use
fixed reciprocal-family
weights for their individual labels.  The incoming power cushions make
their initial excesses vanish; normalize the diagonal-square accounts
with a slightly larger power than their incoming losses.  The initial
normalized total is $n^{o(1)}$.  The estimates above give bounded discount
products for these accounts, so the common
discounted sum has the same extraction margins throughout Phase B.
At a state satisfying the caps, \eqref{eq:original-45} first makes all
candidate precisions positive and all candidate accounts defined.
The minimizing pair then preserves the common sum, and extraction
gives the caps at the next chosen state.  This closes the induction
from $l_2$ to $l_3$.

\subsection{The final matching}

There are now few enough unresolved stubs that the entry bound implies
an operator bound on the remaining boundary error block.  This allows any final
matching.  A separate entry calculation will retain the scale $B_B$
without a factor $l_3$.

Pair the last $l_3$ stubs by any fixed perfect matching.  In this
subsection $P_f$ denotes the involution of this remaining matching,
restricted to the current contrast space; let $V$ be the isometric
inclusion of that space into vertex $\mathbf1^\perp$.
The remaining precision increment is $V(hI-\sigma P_f)V^T$.
Its ideal coefficient is $H_{f0}=-\sigma hP_f$, and the actual
and reference coefficients are
\[
 H_f=H_{f0}(I+DH_{f0})^{-1},\qquad
 \widehat H_f=H_{f0}(I+\widehat D H_{f0})^{-1}.
\]
The entry barrier gives
\[
 \|D\|\le\|\widehat D\|+\|D-\widehat D\|
       \le C/R_0+Cl_3B_B=o(1).
\]
Thus both coefficients are bounded.  The ideal inverse middle
matrix is $-\sigma P_f/h$.  Its eigenvalues have fixed distance
from zero, and its inertia agrees with that of
$m(hI-\sigma P_f)$.  Perturbation by $D=o(1)$ preserves that
inertia.  The Schur-complement inertia identity therefore proves
positivity of the final precision, at every primary point and
both signs.

To retain the entry scale, we use the sparse rows of the reference
inside the exact update.
Let $C_{\rm err}=V^TEV=D-\widehat D$, let
$\varepsilon_i=V^TEe_i$, $p_i=V^T\widehat Pe_i$, and put
$z_i=\widehat H_fp_i$.
The deterministic reference bounds imply
\[
 z_i=\widetilde c_i\Pi_{l_3}e_{t(i)}+\widetilde\tau_i,
 \qquad |\widetilde c_i|\le C,\qquad
 \|\widetilde\tau_i\|\le Cn^{-r_P},
\]
with the coordinate summand possibly absent.  Indeed
$H_{f0}$ sends a short coordinate to its partner, and
$\widehat H_f-H_{f0}$ acting on that coordinate has power-small
norm by the small reference rows and the bounded inverse.
The old entry error yields
\[
 \begin{gathered}
 \|\varepsilon_i\|_\infty\le CB_B,\qquad \|\varepsilon_i\|\le C\sqrt{l_3}B_B,\qquad
 \|C_{\rm err}\|_{\max}\le CB_B,\\
 \max_s\|C_{\rm err}e_s\|\le C\sqrt{l_3}B_B,\qquad
 \|C_{\rm err}\|\le Cl_3B_B.
 \end{gathered}
\]
The exact two-reference Woodbury identity is
\[
 E'_{ij}=E_{ij}-\varepsilon_i^Tz_j-z_i^T\varepsilon_j+z_i^TC_{\rm err}z_j
          -(\varepsilon_i-C_{\rm err}z_i)^TH_f(\varepsilon_j-C_{\rm err}z_j).
\]
It follows by substituting
$H_f-\widehat H_f=-H_fC_{\rm err}\widehat H_f$ in the two update
formulas; thus it does not require either reference matrix to be
invertible.
Every linear displayed term is $O(B_B)$: use its bounded single
coordinate part and the estimate
$\sqrt{l_3}n^{-r_P}=o(1)$ for its tail part.
The two-coordinate term is bounded in the same way, using column
norms for a coordinate--tail term and the operator norm for a
tail--tail term.  In addition
$\|\varepsilon_i-C_{\rm err}z_i\|\le C\sqrt{l_3}B_B$.
The final quadratic term is consequently $O(l_3B_B^2)=o(B_B)$.
This proves the final entry bound in
Proposition~\ref{prop:boundary-completion}.

The completed graph is simple and $d$-regular.  Since no stubs
remain, its precision is
$M(h)=z(h)I-\sigma A(G)$ on $\mathbf1^\perp$.
At the closest point $z(h)<2\sqrt q+e/20$, so positivity gives
\[
 \max_{i\ge2}|\lambda_i(G)|<2\sqrt q+e/20.
\]
The fixed stop configuration and completed matching define the same
reference $\widehat P_{\rm fin}(h)$ throughout the primary interval.
Section~\ref{sec:exact-repair} uses its diagonal derivative bound and
$B_B=o(\sqrt e)$ to obtain local spectral mass bounds.  The retained
stop-path estimates and the aggregate reference formula give the
typical-edge bulk blocks, with a separate exceptional edge set allowed
for each prescribed pair of bulk cutoffs.
Finally, this stage has at most $dn/2$ pair steps, polynomially
many candidate pairs and output labels, a logarithmic number of
retained points, and only fixed-degree moments and finite matrix
operations.  The smooth reference has a fixed finite coefficient
cutoff of length $O(\log n)$, and its coefficients and accounts use
the rational choices specified for the algorithm.  The common
arithmetic-precision argument therefore gives polynomial bit
complexity for all operations in this section.  This completes
the proof of Proposition~\ref{prop:boundary-completion}.
